\documentclass[conference]{IEEEtran}
\usepackage[utf8]{inputenc}
\usepackage{graphicx}

\title{Comparação de Algoritmos de Categorização \\ \large Categorizar um veículo a partir das suas informações \\
\textit{Comparation of Categorization Algorythms \\ \large Categorize a vehicle based on its information}}
\author{
\IEEEauthorblockN{Martinho Caeiro - 23917 || Paulo Abade - 23919}
\IEEEauthorblockA{
    Instituto Politécnico de Beja\\
    Escola Superior de Tecnologia e Gestão\\
    Beja, Portugal\\
    23917@stu.ipbeja.pt || 23919@stu.ipbeja.pt
}
}

\begin{document}
\maketitle

\begin{abstract}
Este artigo apresenta um estudo para a comparação entre algoritmos de categorização de veículos. O objetivo de cada um dos 
algoritmos é categorizar um veículo a partir das suas informações, sendo que este veículo será categorizado consoante o seu 
país de origem. Foi escolhido este tema para facilitar a nossa compreensão sobre o assunto e tornar mais agradável o estudo
destes algoritmos. O estudo foi realizado com base em algoritmos de aprendizagem supervisionada, nomeadamente o algoritmo
\textit{Binary Tree}, algoritmo \textit{Random Forest} TODO: Escolher mais algoritmos para fazer a comparação entre eles. 
Para fazer a comparação destes, foi utilizado um dataset, com cerca de 400 entradas, onde possui informações de veículos de 
diferentes países e este foi utilizado nos diferentes algoritmos para treinar e testar os mesmos. Este dataset, já foi alterado 
no primeiro trabalho da Disciplina de \textit{Sistemas de Informação}, sendo ligeiramente diferente, por já ter sido tratado. 
Os resultados obtidos foram comparados e analisados para perceber qual o algoritmo que melhor categoriza um veículo a partir das 
suas informações. A comparação dos algoritmos foi feita com base na sua precisão, sensibilidade e especificidade. TODO: Adicionar 
mais informações sobre o estudo realizado. Gráficos ROC?
\end{abstract}

\begin{IEEEkeywords}
algoritmos; veículos; categorização; aprendizagem supervisionada; árvore binária; random forest; precisão; sensibilidade; especificidade;
orange; datamining; machine learning; kagoo.
\end{IEEEkeywords}

\section{Introdução}
Este modelo fornece aos autores as principais especificações de formatação necessárias para preparar versões electrónicas dos 
seus artigos obedecendo ao modelo IEEE. Todos os componentes padrão do artigo foram especificados por três razões: (1) facilidade 
de uso para formatação de trabalhos individuais a serem incluídos em atas de conferências, (2) o cumprimento automático de 
requisitos electrónicos que facilitam a produção posterior de produtos electrónicos tais como atas de conferências digitais, 
e (3) a conformidade de estilo ao longo de um livro de atas de conferência.

\section{Facilidade de Utilização}

\subsection{Seleção do Template}
Em primeiro lugar confirme que tem o Template correto, formatado para o tamanho correto de papel (A4).

\subsection{Mantenha as Especificações Corretas}
O modelo é usado para formatar o seu artigo e definir os estilos do texto. Todas as margens, a largura das colunas, espaçamentos 
de linha e fontes de texto são definidos e não devem ser alterados em nenhum caso.

\section{Formatação Geral}
Nos artigos escritos em Português ou Espanhol coloque também título, resumo e palavras-chave em Inglês como indicado.

\subsection{Abreviaturas e Acrónimos}
Definir as abreviaturas e siglas a primeira vez que são utilizadas no texto, mesmo depois de terem sido definidas no resumo. Não 
use abreviatura no título ou títulos de secções a menos que sejam totalmente inevitáveis.

\subsection{Unidades}
\begin{itemize}
    \item Utilize unidades do Sistema Internacional (SI);
    \item Evite combinar unidades SI com outros sistemas (exemplo: corrente em amperes e campo magnético em oersteds) pois gera confusão;
    \item Use zero antes do ponto decimal: “0.25” e não “.25”.
\end{itemize}

\section{Equações}
As equações são uma exceção para as especificações deste modelo. Deve utilizar na equação as fontes Times New Roman e/ou Symbol. 
Numere as equações consecutivamente utilizando parênteses encostados à direita. Exemplo:

\begin{equation}
\alpha + \beta = \chi.
\end{equation}

\section{Figuras e Tabelas}
\subsection{Posicionamento}
Coloque as legendas das figuras por baixo das figuras e os títulos das tabelas em cima das tabelas. Insira as figuras e 
tabelas após serem citadas no texto. Utilize a abreviatura “Fig. 1” para se referir a figuras.

\begin{figure}[h!]
    \centering
    \includegraphics[width=0.5\textwidth]{Resources/Logo/IPBejaESTIG.jpg}
\end{figure}
\newpage
%----------------------------------------------------------------------------------------------------------------------------------------------
\begin{table}[h!]
\caption{Título de uma Tabela Exemplo}
\centering
\begin{tabular}{|c|c|c|}
\hline
Cabeçalho de Coluna & Coluna 1 & Coluna 2 \\
\hline
Linha 1 & Conteúdo 1 & Conteúdo 2 \\
\hline
Linha 2 & Conteúdo 3 & Conteúdo 4 \\
\hline
\end{tabular}
\end{table}

\section{Conclusões}
O modelo IEEE e outros modelos disponíveis facilitam o trabalho de edição de atas de conferências e a produção de material 
científico. Utilize-o corretamente na escrita do seu artigo.

\section*{Agradecimentos}
Introduza agradecimentos a entidades e pessoas que contribuíram para o trabalho. Caso a submissão para a conferência seja 
\textit{Double-Blind}, introduza unicamente o texto \textbf{DFBR} (\textit{Deleted for Blind Review}).

\begin{thebibliography}{1}
\bibitem{ref1} G. Eason, B. Noble, and I. N. Sneddon, “On certain integrals of Lipschitz-Hankel type involving products 
of Bessel functions,” \textit{Phil. Trans. Roy. Soc. London}, vol. A247, pp. 529–551, April 1955.
\end{thebibliography}

\end{document}
